%%%%%%%%%%%%%%%%%%%%%%%%%%%%%%%%%%%%%%%%%%%%%%%%%%%%%%%%%%%%%%%%%%%%%%%%%%%%%%%%
%%
\cleardoubleoddpage%  Make sure to start each chapter on a new odd page
\chapter{Discussion and Conclusion}
\label{sec:conclusion}

\section{Discussion}

The reproduction succeeded in the sense that matters for this practical
work: not bit-exact agreement with the paper's plotted curves, which is not
achievable given the TensorFlow/Keras to PyTorch/JAX framework
substitution that predates this report, but agreement in the paper's
qualitative claims, using the paper authors' own mesh and finite element
discretization rather than an independently calibrated approximation of
one (Appendix~\ref{app:an-appendix}). Across both the Hilbert-valued
diffusion experiment and the Banach-valued NSB experiment, the reproduced
error decays at approximately the paper's reported $O(m^{-1})$ rate (mean
fitted slopes $-1.02$ for diffusion and $-0.99$/$-0.96$ for NSB $u$/$p$),
and ELU beats ReLU by a wide, consistent margin everywhere, the single
most robust finding in the whole study. tanh, however, does not beat ReLU
uniformly the way the paper's blanket "ELU/tanh beat ReLU" claim suggests.
tanh matches ELU closely at the smaller (4$\times$40) architecture size,
but is consistently the worst performer at the larger (10$\times$100)
size, in both experiments, both coefficient families, and both
dimensions. So what actually shows up is an interaction between activation
and architecture size, not a flat per-activation ranking
(Section~\ref{sec:experimentalresults}). The claim of no degradation from
$d=4$ to $d=8$ is the other place the reproduction diverges from a clean
match: dimension robustness tracks architecture size in both experiments
(small networks close to neutral, large networks degrading, large tanh
worst of all), which is one consistent driver across both PDEs, unlike an
earlier, uncorrected version of this reproduction's pipeline, which had
suggested two different drivers. Both findings are reported as genuine
empirical properties of the paper authors' own discretization, run
exactly as they implemented it, rather than reconciled away, since they
are consistent across coefficient families within each PDE.

PyTorch and JAX agree with each other closely under the identical protocol
described in Section~\ref{sec:jaximplementation}: relative test error
differs by well under 6.5\% between frameworks in every one of the twelve
cases (Section~\ref{sec:framework-comparison-results}). This close
agreement is itself evidence that any remaining mismatch with the paper's
own reported numbers comes from the framework substitution and the one
remaining documented deviation, NSB's reduced training budget
(Appendix~\ref{app:an-appendix}), rather than from an error specific to
this reproduction's own pipeline. If the pipeline itself were unreliable,
two independently implemented frameworks trained on identical data would
have no particular reason to agree this closely. Training time tells a
more structured story than either framework's raw speed would suggest. For
diffusion, PyTorch is consistently faster than JAX in every case
(19--25\% less time). For NSB the split is not case- but target-dependent:
PyTorch is faster for velocity $u$ in all four cases, JAX is faster for
pressure $p$ in all four cases, with no exceptions in either direction.
This is the opposite ranking, for diffusion, from what
Section~\ref{sec:jaximplementation} anticipated (JIT compilation
amortizing better for JAX on the larger, more numerous diffusion runs).
The clean $u$/$p$ split for NSB, where both targets share identical
architectures and shape diversity, suggests output dimension itself (1464
DOF for $u$ vs.\ 244 for $p$) is the more likely driver rather than
compiled-shape diversity, though this reproduction does not have a
confirmed mechanism for it (Section~\ref{sec:framework-comparison-results}).

\section{Limitations}

The stationary Boussinesq equation, the paper's third experiment, was out
of scope for this practical work (Section~\ref{sec:problemStatement}). The
diffusion mesh, the NSB finite element family, the NSB inlet velocity, and
the NSB pressure gauge are no longer independently reconstructed or
calibrated approximations. They are read directly from the paper authors'
own released implementation and verified against it exactly
(Appendix~\ref{app:an-appendix}), which removes what was previously this
report's largest source of uncertainty. Two genuine deviations remain. The
sparse-grid test level is not stated numerically in the paper's own text,
but is now taken directly from the authors' own batch scripts rather than
chosen independently. NSB training uses a reduced epoch budget and a
looser loss tolerance than the paper specifies
(Appendix~\ref{app:an-appendix}), a deviation made necessary by measured
wall-clock cost rather than by any property of the paper's method itself.
Section~\ref{sec:experimentalresults} shows this most plausibly explains
NSB's slightly shallower average fitted decay rate relative to diffusion,
and it means the NSB numbers in this report are not run under a strictly
paper-identical training budget the way the diffusion numbers are.
Finally, all absolute error magnitudes in this report are specific to the
PyTorch/JAX implementations described above; they should not be read as
literal reproductions of the paper's own plotted values, only as a check
on the paper's qualitative claims.

\section{Conclusion}

Both PyTorch and JAX reproductions of the diffusion and NSB experiments
recover the paper's central qualitative findings, an approximately
$m^{-1}$ error-decay rate and a clear (if size-dependent)
activation-function ranking, using the paper authors' own mixed finite
element mesh, element discretization, and sparse-grid test protocol
rather than an independently reconstructed stand-in, at the paper's full
experimental scale (2 PDEs, 2 coefficient families, 2 dimensions, 6
architectures, 14 training sizes, 12 trials, both frameworks: 24,192
completed training runs, 8,064 diffusion and 16,128 NSB, with zero
failures).

Implementing the paper's protocol precisely rather than approximately
turned out to matter in several concrete ways that a looser reproduction
would have missed entirely. The diffusion mixed formulation's natural
(rather than essential) boundary condition changes what "implementing the
boundary condition" even means at the code level. Keras'
activation-agnostic \texttt{HeUniform} initializer is a silent, easy-to-miss
mismatch against both PyTorch's and JAX's own activation-aware defaults.
The paper's checkpoint/restore rule, read directly from the authors' own
training callback, turned out to be a single best-loss-so-far trigger, not
the two-trigger rule an initial reading of the paper's prose appendix
suggested. And the paper's test-error formula, a FEM mass-matrix-weighted
Bochner norm against sparse-grid quadrature weights that are genuinely
negative at a real fraction of points, silently produces NaN under a
naive implementation of the same formula unless each accumulated weighted
sum is wrapped in an absolute value before the square root, exactly as the
authors' own code does. Each of these was only caught by implementing the
letter of the paper's appendix and cross-checking directly against the
authors' own released code, rather than a plausible approximation of
either. Recovering the paper's qualitative claims despite these
implementation-level departures from a naive reading of the method
description suggests the paper's reported behavior is a robust property of
the method, not an artifact of one specific framework or one specific
finite element choice. At the same time, it also revealed that the
paper's own "ELU/tanh beat ReLU, no degradation at $d=8$" claims hold only
for a subset of the architectures the paper itself tests, not universally.

The most natural next step for extending this reproduction is adding the
stationary Boussinesq equation, the paper's third experiment, which was
excluded from this practical work's scope
(Section~\ref{sec:problemStatement}) but would complete the paper's full
experimental picture using the same infrastructure built here. A second
natural extension is re-running the NSB experiment under the paper's full,
un-reduced training budget (Appendix~\ref{app:an-appendix}) given more
compute time than was available for this practical work, to check whether
NSB's slightly shallower fitted decay rate closes the small remaining gap
to diffusion's. A third is investigating why large (10$\times$100) tanh
networks specifically fail to fit the corrected, higher-fidelity
discretization while every other architecture/activation combination
converges cleanly. This reproduction establishes the pattern empirically
(Section~\ref{sec:experimentalresults}) but does not diagnose its cause,
for example whether it comes from tanh saturation interacting with a
wider network's larger pre-activation values, or from an
optimization/learning-rate-schedule mismatch specific to that architecture
size.

%%
%%%%%%%%%%%%%%%%%%%%%%%%%%%%%%%%%%%%%%%%%%%%%%%%%%%%%%%%%%%%%%%%%%%%%%%%%%%%%%%%
